\begin{textsample}{POS Dim 4 – ce – Score 8.00 – ce/ce\_inf\_e\_ef\_28.txt}  \label{ex:f4_pos_139}
3.12.
Relações de Gênero A Educação Infantil é a primeira etapa da educação básica e tem como finalidade o desenvolvimento integral da criança de 0 a 5 anos de idade, em seus aspectos físicos, psicológicos, intelectuais e sociais, complementando a ação da família e da comunidade.
O Ensino Fundamental, etapa subsequente, tem como objetivos, conforme Artigo 32 da LDB : “ o desenvolvimento da capacidade de aprendizagem, tendo em vista a aquisição de conhecimentos e habilidades e a formação de atitudes e valores ; o fortalecimento dos vínculos de família, dos laços de solidariedade humana e de tolerância recíproca em que se assenta a vida social ” ( BRASIL, 1996 ).
Partindo desses pressupostos legais, compreendemos que temos na sociedade meninos e meninas com histórias de vida diversificadas.
Eles nos motivam a pensar o quão é essencial que a escola incorpore e discuta democraticamente as questões que envolvem as relações entre sujeitos, considerando o respeito e a autonomia como basilares ao relacionamento humano.
Isso possibilita uma ampla crítica cultural do domínio masculino e da subordinação feminina baseada em relações desiguais de gênero.
Biologicamente falando, meninos e meninas não agem de maneira diferente em consequência de sua fisiologia, mas como decorrência de uma educação culturalmente apropriada e diferenciada a partir do sexo.
São educados com base em convenções de adultos que tratam as crianças como meros figurantes e assim vão definindo o que julgam importante para a vida dessas pessoas, a exemplo das brincadeiras e/ou das cores que devem usar.
A ideia de que as meninas devem brincar de certos jogos e brinquedos ( bonecas, panelas... ) e os meninos, com outros ( carros, bolas, armas... ) tem como fundamento o patriarcado[^2 ] e machismo[^3 ], geralmente velados nas práticas educativas.
A lógica de associar a recreação e o brincar das meninas a objetos e brincadeiras que remetem ao mundo do cuidado e afetos ( casinha, bonecas etc. ) e, da mesma forma, atribuir aos meninos objetos e brincadeiras associadas a um universo público em detrimento do âmbito privado reflete a sociedade sexista em que vivemos.
Contribuir com essa perpetuação é limitar o processo de aprendizagem nas atividades educativas e estimular preconceitos, \textbf{discriminações}, misoginia,[^4 ] racismos e \textbf{intolerância} em vários níveis.
Tudo isso corrobora com os altos índices de violência e assassinatos de meninas, \textbf{mulheres} e demais sujeitos sociais.
Cabe à escola instigar o debate acerca do respeito que deve permear as relações sociais e todas as dimensões que o envolvem.
Por isso é que construímos a escola no espaço relevante para a reconstrução de uma nova ordem social, igualitária, humana e fraterna.
Julgamos necessário rever as orientações aparentemente inofensivas às crianças, perpetradas por uma cultura segregadora, uma vez que isso se reflete nos resultados delas e deles e em como deverão enxergar o mundo, quando na vida adulta.
É exatamente em função de um tratamento diferenciado entre meninas e meninos na escola e na sociedade em geral que percebemos os nichos profissionais ocupados por um e/ou por outro grupo.
Aos \textbf{homens} são atribuídas funções cujo perfil é a racionalidade ; às \textbf{mulheres}, as profissões que exigem o comportamento materno e o cuidar.
Consequentemente, a disparidade salarial entre esses constructos sociais é, também, absurda, embora concebida como natural.
Faz -se \textbf{urgente} a necessidade de reformularmos alguns hábitos, conceitos e preconceitos naturalizados em nossa sociedade.
Um dos caminhos para a construção dessa sociedade justa e com equidade é o trabalho que a escola pode desenvolver desde a Educação Infantil.
A escola precisa assumir o seu papel na construção do desenvolvimento integral como está garantido às crianças na LDB.[^5 ] Por isso, é essencial que a instituição mantenha diálogo com as famílias sobre a importância do respeito nas relações interpessoais, da autonomia e, sobretudo, que exponha uma abordagem não sexista do conteúdo didático.
Educadoras/es precisam compreender que, durante a infância e a adolescência, as pessoas estão em processo de formação em sua \textbf{identidade}, seus desejos e vontades subjetivas.
Determinar comportamentos, brincadeiras, objetos, cores, roupas “ para meninos ” e “ para meninas ” gera segregação sexista na compreensão da realidade das crianças.
É nesse tipo de comportamento que o significado do lugar social se solidifica numa perspectiva de predeterminação.
Esse pensamento e ação presentes nas escolas causam distorções muito graves na construção do sujeito.
Eles propagam uma concepção de mundo que \textbf{hierarquiza} pessoas e suas capacidades, a exemplo da ideia \textbf{cristalizada} de que as \textbf{mulheres} devem receber salários inferiores, ou devem ocupar postos de trabalho relacionados ao cuidado.
Para além disso, observamos a objetificação do corpo e da vida da \textbf{mulher}, tornando comum o assédio sexual, como se fosse um direito do \textbf{homem}.
Precisamos considerar, ainda, sobre essa questão, certa confusão entre categorias do pensamento social humano : gênero, sexo e sexualidade.
Embora digam respeito à subjetividade e à singularidade do ser humano ( no sentindo ontológico ), detêm significados diversos, mas, em função de um discurso ideológico, meninas são tratadas como inferiores em relação aos meninos.
A negação de direitos como jogar futebol, correr, soltar pipas, brincar com carrinho, dentre outros, é uma realidade no cotidiano de muitas escolas.
Em contrapartida, os meninos não têm seus corpos e brincadeiras vigiados da mesma forma.
Tal posicionamento ratifica a ideia de que o \textbf{homem} “ pode tudo ” : correr, ser livre, falar alto, inclusive usar brincadeiras agressivas ( lutas, armas ).
Assim, vigia -se o andar, as brincadeiras, o contato com as cores, evitando ao máximo uma aproximação de meninos no que diz respeito aos afetos.
Quando nos deparamos com meninos que, durante o horário do lanche, se recusam a colocar o alimento em recipientes de cor rosa, por exemplo, percebemos como a violência de gênero é capaz de modificar o comportamento das pessoas através de práticas banais.
Isso exige de nós mudança de postura e análise sobre uma sociedade eivada de desrespeito, \textbf{discriminação} e preconceito.
O direito a uma educação que preza pela qualidade inclui a discussão e análises das questões de \textbf{identidade}, respeito e autonomia junto ao corpo docente.
As relações entre as crianças, nesta etapa da educação, apresentam -se como uma das primeiras formas de introdução de meninos e meninas na vida social.
Isso decorre porque possibilitam a oportunidade de terem contato com crianças das mais variadas classes sociais, religiões, culturas e etnias, com comportamentos e valores também diferenciados.
Cabe lembrar que a educação reflete um modelo de sociedade.
Se ele é desigual, certamente contribuirá para a manutenção de uma dinâmica social que define lugares e papeis, inclusive no âmbito dos afetos.
Organizada nesta perspectiva, a sociedade produz lacunas entre os sujeitos, provocando sofrimento, tanto para os que sentem a opressão no cotidiano, quanto para o sujeito que, mesmo sem reconhecer, é o agente ativo da opressão.
Numa sociedade em que as \textbf{mulheres} assumem o espaço do cuidado com a prole e o acompanhamento do rendimento escolar de filhos, enquanto a maioria dos \textbf{homens} é afastada dessa vivência de afetos para responder a um padrão de masculinidade, existe uma relação de desigualdade e, principalmente, de opressão.
Nesse contexto ideológico, a educação precisa atuar de forma a romper com a lógica da desigualdade, elevando -se a um nível que considere os sujeitos sociais de forma inclusiva, respeitosa e humana.
A escola não pode somente cuidar dos aspectos intelectuais, da higienização do corpo ou do ensino dos códigos alfanuméricos ; ela precisa também considerar as várias linguagens e os diversos aspectos históricos e sociais presentes nas salas de aula e na vida de educadoras/es e educandas/os.
Evidente que observar aspectos tão negligenciados no processo educativo, mesmo que elementares, só é possível a partir de um modelo de educação que considere o contexto educacional de alunas e alunos.
Devemos praticar um processo educativo que identifica a/o estudante como protagonista e considera, além da relação com o meio ambiente, as relações sociais ali estabelecidas.
Portanto, \textbf{identidades}, relações sociais de desigualdade e enfrentamento à violência contra a \textbf{mulher} são questões que podem ser consideradas didaticamente pelas várias disciplinas da Educação Infantil e Ensino Fundamental, citando apenas dois níveis de educação conforme Artigo 21 da LDB 9394/96.
Nesse sentido, necessitamos urgentemente que a escola assuma o debate em torno da temática \textbf{identidade}, do respeito e da autonomia no âmbito das suas atividades cotidianas que ultrapassem os limites dos cuidados com as estatísticas exigidas nas avaliações.
É \textbf{urgente} a efetiva ressignificação das relações humanas e sociais, acenando as possibilidades da construção de uma sociedade mais justa, inclusiva, democrática, participativa e não violenta.
Pensar nas questões do respeito às relações interpessoais e da autonomia na Educação Infantil e Ensino Fundamental, para além das discussões dicotômicas entre masculino e feminino, é pensar em ações que busquem garantir a equidade, pois diz Santos ( 2003, p. 53 ) : “ Lutar pela igualdade sempre que as diferenças nos discriminem ; lutar pelas diferenças sempre que a igualdade nos descaracterize. ” [ ^2 ] : Patriarcado : “ O patriarcado refere -se a milênios da história mais próxima, nos quais se implantou uma hierarquia entre \textbf{homens} e \textbf{mulheres}, com primazia masculina...
É a esta estrutura de poder, e não apenas à ideologia que a acoberta, que o conceito de patriarcado diz respeito ” – Heleieth Saffioti. [ ^3 ] : Machismo : é a base de sustentação, produto do patriarcado, que é o sistema que, na construção do \textbf{Estado}, colocou as \textbf{mulheres} na esfera privada ( Agencia Brasil ). [ ^4 ] : Misoginia : ação de repulsa ou aversão às \textbf{mulheres}, aversão excessiva do contato sexual com \textbf{mulheres}. ( Etm.
Do grego misos : ódio e gene : \textbf{mulher} ).
Dicionário Online. [ ^5 ] : LDB — Lei de Diretrizes e Bases da Educação

% matched lemmas: cristalizar, discriminação, estado, hierarquizar, homem, identidade, intolerância, mulher, urgente
\end{textsample}
